\documentclass[10pt,a4paper]{amsart}
\usepackage[margin=19mm]{geometry}
\usepackage{amsmath,amssymb,mathtools}
\usepackage[scr=rsfs,cal=euler]{mathalfa}
\usepackage{xfrac}
\usepackage[unicode,hidelinks,bookmarks=false]{hyperref}
\newcommand{\ie}{i.e.\ }
\newcommand{\cf}{cf.\ }
\newcommand{\resp}{respectively\ }
\newcommand{\Subsection}{\subsection}
\providecommand{\nobreakdash}{\nobreak\hbox{-}}
\setlength{\emergencystretch}{2em}
\multlinegap0pt
\usepackage{amssymb}
\usepackage{enumerate}
\usepackage{xcolor}
\newtheorem{theorem}{Theorem}
\newcommand{\bc}{\mathbb{B}}
\newcommand{\dbar}{\bar\partial}
\title{Generalizations of the Dirichlet problem for bianalytic functions}
\author{William L. Blair}
\date{Source: arXiv:2605.18573v1 (CC BY 4.0)}
\begin{document}
\maketitle

\noindent\emph{Source and licence.} Generalizations of the Dirichlet problem for bianalytic functions. Version arXiv:2605.18573v1 (CC BY 4.0), distributed by arXiv under the Creative Commons Attribution 4.0 International licence, http://creativecommons.org/licenses/by/4.0/, as recorded in the arXiv licence metadata for this exact version and shown on the abstract page https://arxiv.org/abs/2605.18573v1. Full licence notice: \texttt{licenses/arxiv-2605.18573-CC-BY-4.0.txt}.\par\smallskip

\noindent\textit{Editorial scope.} Counted unit: the article's theorem bchoivsimprinreptheorem (source0.tex lines 902--912). The article's complete original proof of it is delivered with it (source0.tex lines 1926--1944, one \texttt{proof} environment of 19 lines, which the article places away from the statement). No mathematics is written, repaired or re-derived: the statement and the proof are the article's own lines, in the article's own notation, and no retained line is substituted or re-flowed.  Retained uncounted as the source-local context the proof needs: lines 1375--1384.  Nothing was omitted from any retained span, so the package carries no omission record.  No retained line contains a citation, so the wrapper carries no bibliography and no bibliography entry.  No retained line contains a figure environment, an image-inclusion command or drawing code, so no image is delivered and the package declares no asset dependency.  Editorial formatting only, in addition: an AMS \texttt{amsart} preamble that restates the article's own declarations and macros used by the retained lines, namely the environments \texttt{theorem} on separate counters, together with the standard packages those lines need.  The retained environments print in this document as Theorem 1 in order of appearance, whereas the article prints the same items with the numbering of its own full document.  The complete native source of the article as distributed in the arXiv e-print for this exact version is bundled unchanged as \texttt{sources/200865/source0.tex}.
            \begin{theorem}\label{thm: generalizationOfBianalyticWellPosed}
             Let $Q(x,y) = 0$ be a non-degenerate conic that is not a circumference, $\Gamma = \{(x,y) \in \mathbb{R}^2 : Q(x,y) = 0\}$, and $P$ be a polynomial in $x$ and $y$. The Dirichlet problem for the $\bc$-Bitsadze equation
                \begin{equation}\label{eq: bicomplexdirichletthm}
                    \begin{cases}
                        \dbar^2 w  = 0, & \text{ in } \mathbb{R}^2\setminus \Gamma,\\
                        w = P, & \text{ on } \Gamma,
                    \end{cases}
                \end{equation}
                has a unique solution.
                \end{theorem}

            \begin{theorem}\label{thm: bchoivsimprinreptheorem}
                For a bounded simply connected set $D \subset \mathbb{C}$, let $n$ be a positive integer and $A \in W^{n-1, q}(D,\bc)$, $q>2$. Every solution $w: D \to\bc$ of 
                \[
                    (\dbar- A)^n w = 0
                \]
                is representable as 
                \[
                    w(z) = e^{T_D^\bc[A](z)}\sum_{k=0}^{n-1} \widehat{z^*}^k h_k(z),
                \]
                where $h_k\in Hol(D,\bc)$, for every $k$.
            \end{theorem}

            \begin{proof}[Alternative Proof]
                All solutions of $(\dbar - A)^2 w = 0$ have the form 
                \[
                    w = e^{T_D^\bc[A]} (\varphi_0 + \widehat{z^*}\varphi_1),
                \]
                where $\dbar \varphi_i = A\varphi_i$, $i \in \{1,2\}$, by Theorem \ref{thm: bchoivsimprinreptheorem}. By Theorem \ref{thm: generalizationOfBianalyticWellPosed}, there exist unique $\varphi_0, \varphi_1$ such that $b = \varphi_0 + \widehat{z^*} \varphi_1$ solves
                \[
                    \begin{cases}
                        \dbar^2 b = 0 & \text{ in } D\\
                        b = P & \text{ on } \Gamma
                    \end{cases}.
                \]
                Therefore, $w = e^{T_D^\bc[A]} b$ solves $(\dbar - A)^2 w = 0$ and 
                \begin{align*}
                    w|_{\Gamma} &= e^{T_D^\bc[A]} b|_{\Gamma}
                    = e^{T_D^\bc[A]} P. 
                \end{align*}
                Uniqueness of the solution follows from the uniqueness of $b$. 
            \end{proof}

\end{document}
